\documentclass[12pt]{article}
\usepackage[a4paper,margin=2.2cm]{geometry}
\usepackage{amsmath,amssymb,amsthm,mathtools}
\usepackage{hyperref}
\usepackage{cleveref}

% 中文（xelatex）
\usepackage{fontspec}
\usepackage{xeCJK}
\setCJKmainfont{Noto Serif CJK SC}

\title{Math-Note- Algebraic enhancement training}
\date{}
\begin{document}
\maketitle

\section{LCD}
\subsection*{LCD 清除法 The LCD Method}
\subsection*{题目1 : Solve for $\frac{2 + \frac{1}{x}}{4 -\frac{1}{x}}$}

\subsection*{思路:}
分母:  $4 - \frac{1}{x}$  存在分数 $\frac{1}{x}$
我们需要先消分?

\subsection*{步骤:}
1:  $\frac{2 + \frac{1}{x}}{4 -\frac{1}{x}} \cdot x $   \\
2:  = $\frac{2x + 1}{4x - 1}$ ? 还能继续化简吗? 

\subsection*{题目2: Solve for  $\frac{\frac{3}{a}-\frac{3}{b}}{b-a}$}

\subsection*{思路:}
在代数化简的思路 我们首先可以先尝试 解决 分数 或 负指数的消除
我们发现 分子是存在 分数 $\frac{3}{a}-\frac{3}{b}$, 它们的LCD是a,b

\subsection*{步骤:}
1:  $\frac{3}{a}-\frac{3}{b} \cdot ab$ \\
2:  因为分子*ab, 分母也需要*ab-> $\frac{3b -3a}{(b-a)\cdot ab}$ \\
3:  提取公约数和调整位置  $\frac{3(b-a)}{ab(b-a)}$ \\
4:  约分  $\frac{3}{ab}$ \\

\subsection*{答案:}
$\frac{3}{ab}$ 

\subsection*{题目3: Solve for  $\frac{(x+h)^{-1} - x^{-1}}{h}$}
\subsection*{步骤:}
1: $(x+h)^{-1} - x^{-1}$ ->  $\frac{1}{(x+h)} - \frac{1}{x}$ \\
那么LDC 是 x(x+h) \\
2 : $(\frac{1}{(x+h)} - \frac{1}{x}) \cdot x(x+h)$ \\ -> $x- (x+ h)$ \\
3 : 因为分子乘以了 x(x+h), 分母也需要乘以 x(x+h) ->  $\frac{x - (x+h)}{h \cdot x(x+h)}$ \\
4 : 简化并约分 $\frac{-h} {h \cdot x(x+h)}$  ->  $\frac{-1}{x \cdot(x+h)}$

\subsection*{答案:}
$\frac{-1}{x \cdot(x+h)}$ 


\section{Refactoring Exponents and Radicals:} 

\subsection*{基础转换 (Basic Casting)}

\subsection*{题目: Conversion  $\sqrt[5]{x}$}
\subsection*{答案:}
$x^\frac{1}{5}$


\subsection*{题目: Conversion  $\frac{1}{x^4}$}
\subsection*{答案:}
$x^{-4}$


\subsection*{题目: Conversion  $\sqrt{x^{3}}$}
\subsection*{答案:}
$x^{\frac{3}{2}}$


\section{Batch Processing}
\subsection*{题目: Conversion  $x^{2}\cdot\sqrt{x}$}
\subsection*{思路:}
1: 先把所有项转换成指数形式 \\
$\sqrt{x}$ ->  $x^{\frac{1}{2}}$  \\
2: 应用指数法则 \\ 
$x^{2} \cdot x^{\frac{1}{2}}$  ->  $x ^ {\frac{1}{2} + 2}$ \\

\subsection*{答案:}
$x^{\frac{5}{2}}$
\subsection*{题目: Conversion  $\frac{x^{3}}{\sqrt{x}}$}
\subsection*{思路:}
1: 先把所有项转换成指数形式 \\
$x^{3}\cdot\frac{1}{x^{\frac{1}{2}}}$ \\
$x^{3}\cdot x^{-\frac{1}{2}}$ \\
$x^{3 - \frac{1}{2}}$ \\
\subsection*{答案:}
$x^{\frac{5}{2}}$ \\


\subsection*{Batch Processing}
\subsection*{题目: Conversion  $\frac{1}{2\sqrt{x}}$}
\subsection*{思路:}
1: 先把所有项转换成指数形式 \\
$\frac{1}{2} \cdot \frac{1}{x^{\frac{1}{2}}}$ \\
$\frac{1}{2} \cdot x^{-\frac{1}{2}}$ \\

\subsection*{答案:}
$\frac{1}{2}x^{-\frac{1}{2}}$

\subsection*{题目: Conversion  $\begin{pmatrix}\frac{2x^{3}}{y^{2}}\end{pmatrix}^{3}$}
1: 解外指数
$\frac{2^{3}x^{3*3}}{y^{2*3}}$
$\frac{8x^{9}}{y^{6}}$
1: 转换分数
$8x^{9}y^{-6}$
\subsection*{答案:}
$8x^{9}y^{-6}$

\subsection*{题目: Conversion  $\sqrt{x\sqrt{x}}$}
1: 解决外根号
 $\begin{pmatrix}
 x\sqrt{x}
\end{pmatrix}^{\frac{1}{2}}$
1: 解决内根号
 $\begin{pmatrix}
 x^{1} \cdot x^{\frac{1}{2}}
\end{pmatrix}^{\frac{1}{2}}$
 $\begin{pmatrix}
 x^{\frac{1}{2} + 1}
\end{pmatrix}^{\frac{1}{2}}$
 $\begin{pmatrix}
 x^{\frac{3}{2}}
\end{pmatrix}^{\frac{1}{2}}$
2: 解决外指数
 $x^{\frac{3}{2} \cdot \frac{1}{2}}$
 $x^{\frac{3}{4}}$

\subsection*{答案:}
= $x^{\frac{3}{4}}$

\section{Polynomial Expansion and Structure}


\subsection*{Basic Squares}
\subsection*{题目: Unfold:  $(x+5)^{2}$}

\subsection*{答案}: \\
= $x^{2} + 2(x5) + 5^{2}$ \\
= $x^{2} + 10x + 25$

\subsection*{题目: Unfold:  $(2x-3)^{2}$}
$(2x)^{2} - 2\cdot2x\cdot3 + 3^{2}$ \\
= $4x^{2} - 12x + 9$ \\
= $4x^{2} - 12x + 15$  \\

\subsection*{答案:} 
= $4x^{2} - 12x + 15$ 

\subsection*{Pre-Calculus Drill}
\subsection*{题目: Unfold:  $f(x) = 3x^{2} + 2$  计算并展开 $f(x+h)$}
\subsection*{步骤:}
1: 展开代元
$3(x+h)^{2} + 2$ \\
= $3(x^{2} + 2xh + h^{2}) + 2$ \\
2: 分配律展开 \\
= $3x^{2} + 6xh + 3h^{2} + 2$ \\

\subsection*{答案:} 
\begin{align*}
&= 3x^{2} + 6xh + 3h^{2} + 2 \\
\end{align*}

\subsection*{The Cube Challenge}
\subsection*{题目: Unfold:  计算并展开 $(x+2)^{3}$}
\subsection*{步骤:}
1: 使用立方展开公式 \\
$x^{3} + 3x^{2}2 + 3x2^{2} + 2^{3}$ \\
2: 计算各项 \\
= $x^{3} + 6x^{2} + 12x + 8$ \\
\subsection*{答案:}
\begin{align*}
    & = x^{3} +6x^{2} + 12x + 8 \\
\end{align*}

\subsection*{题目: Unfold:  计算并展开 $\frac{(x+h)^{3}-x^{3}}{h}$}
1: 使用立方展开公式 \\
\begin{align*}
&= \frac{x^{3} + 3x^{2}h  + 3xh^{2} + h^3 - x^{3}}{h}
&= \frac{3x^{2}h  + 3xh^{2} + h^3}{h} \\ 
&= \frac{h(3x^{2} + 3xh  + h^{2})}{h}
& = 3x^{2} + 3xh  + h^{2} \\
\end{align*}
\subsection*{答案:}
\begin{align*}
   & = 3x^{2} + 3xh  + h^{2} \\
\end{align*}


\subsection*{题目: Unfold:  计算并展开 $(2x-3)^{2}$}
= $(2x)^{2} - 2\cdot2x\cdot3  + 3^{2}$  \\
= $4x^{2} - 12x  + 9$

\subsection*{答案:}
$4x^{2} - 12x  + 9$


\subsection*{题目: Unfold:  计算并展开 $(3x-5)^{2}$}
= $3x^{2} - 2\cdot3x\cdot5 - 5^{2}$
= $9x^{2} - 30x + 25$
\subsection*{答案:}
= $9x^{2} - 30x + 25$

\subsection*{题目: Unfold:  展开 $(3x-5)\cdot(3x-5)$}
\begin{align*}
  &= 3x * 3x  + 3x * (-5)  + (-5) * 3x + (-5) * (-5) \\
  &= 9x^{2} - 15x - 15x  +25 
  &= 9x^{2} - 30x + 25
\end{align*}

\section{Factoring: Reverse Engineering}

\subsection*{Pattern Matching}
\subsection*{题目: 分解 Factor: $4x^{3} - 8x^{2}$}
\subsection*{答案:}
\begin{align*}
  &= 4x^{2}(x - 2) \\
\end{align*}

\subsection*{题目: 平方差 Factor: $x^{2} - 36$}
\subsection*{答案:}
\begin{align*}
  &= x^{2} - 6^{2} \\
  &= (x-6)(x+6) \\
\end{align*}

\subsection*{题目: 二次三项式 Factor: $x^{2} + 7x + 10$}
\subsection*{答案:}
\begin{align*}
  &= (x+2)(x+5)\\
\end{align*}

\subsection*{Combined Logic}
\subsection*{题目: 分解 Factor $2x^{2} - 8$}
\subsection*{思路}
= $2x^{2}-2^{3}$\\
= $2^{1}x^{2}-2^{3}$\\
= $2(x^{2}-2^{2})$\\
= $2(x^{2}-4)$\\
$x^{2}-4$ 是否还可以继续分解\\
= $x^{2} + 0x -4$\\
= $(x+2)(x-2)$\\
= $2((x+2)(x-2))$\\
 
\subsection*{题目: 分解 Factor $x^{2}-2x-15$}
\subsection*{思路}
寻找两个数，乘积等于常数项，和等于一次项系数
  b = -15, a = -2
  这两个数是-5,3
\begin{align*}
  & =(x-5)(x+3) \\
  check-> \\
  &= x^{2}+3x-5x-15 \\ 
  &= x^{2}-2x-15 \\
\end{align*}

\subsection*{答案:}
\begin{align*}
  &= (x-5)(x+3) \\
\end{align*}

\subsection*{Calculus Context}
\subsection*{题目: 分解 Factor $\frac{x^{2}-x-6}{x-3}$}

\subsection*{思路}
先处理分子 a= -1, b= -6
这两个数是-3,+2
= $(x-3)(x+2)$\\
= $\frac{(x-3)(x+2)}{x-3}$\\
约掉 x-3
= $x+2$\\

\subsection*{答案:}
\begin{align*}
  &= x+2 \\
\end{align*}

\subsection*{Higher Order}
\subsection*{题目: 分解 Factor $x^{4}-16$}
\subsection*{思路}
= $(x^{2})^{2} - 4^{2}$
设 $y=x^{2}$ 
= $y^{2} - 4^{2}$
使用平方差公式
= $(y-4)(y+4)$
代回来
= $(x^{2}-4)(x^{2} + 4)$ 
= $((x+2)(x-2))(x^{2} + 4) $

\subsection*{题目: 化简以下式子，使其不含负指数： $\frac{6x^{-2}y^{3}}{2x \cdot y^{-2}}$}
\subsection*{思路}
= $\frac{6\frac{1}{x^{2}}y^{3}}{2x \cdot \frac{1}{y^{2}}}$\\
= $\frac{6 \cdot \frac{y^{3}}{x^{2}}}{2 \cdot \frac{x}{y^{2}}}$\\
= $\frac{6}{2} \cdot \frac{\frac{y^{3}}{x^{2}}}{\frac{x}{y^{2}}}$  提出常数\\
= $\frac{6}{2} \cdot (\frac{y^{3}}{x^{2}})÷(\frac{x}{y^{2}})$      分数相除\\
= $\frac{6}{2} \cdot (\frac{y^{3}}{x^{2}})\cdot(\frac{y^{2}}{x})$  乘倒数\\
= $3 \cdot \frac{y^{3} \cdot y^{2}}{x^{2} \cdot x^{1}}$\\
= $3 \cdot \frac{y^{5}}{x^{3}}$\\
= $\frac{3y^{5}}{x^{3}}$\\
\subsection*{答案:}
\begin{align*}
  &= \frac{3y^{5}}{x^{3}} \\
\end{align*}

\subsection*{题目: 化简以下式子，使其不含负指数： $(3x-4)^{2}$}
\subsection*{思路}
= 使用开二次方公式 $(a-b)^{2}$ =  $a^{2}-2ab + b^{2}$\\
= $(3x)^{2} -24x + 16$\\
拆开推\\
= $(3x-4) \cdot (3x-4)$\\
= $3x * 3x - 12x - 12x  + 16$\\
= $(3x)^{2} -24x + 16$\\
\subsection*{答案:}
\begin{align*}
  &= (3x)^{2} -24x + 16 \\
\end{align*}

\subsection*{题目: 化简： $\frac{\frac{1}{(x+h)^{2}} - \frac{1}{x^{2}}}{h}$}
\subsection*{思路}
解决通分
$(x+h)^{2}$ 和 $x^{2}$ 的公分母是 $(x+h)^{2}x^{2}$\\
= $\frac{\frac{x^{2}}{(x+h)^{2}x^{2}} - \frac{(x+h)^{2}}{(x+h)^{2}x^{2}}}{h}$\\
= $\frac{\frac{x^{2} - (x+h)^{2}}{(x+h)^{2}x^{2}}}{h}$\\
= $\frac{\frac{x^{2}-(x^{2} + 2xh + h^{2})}{(x+h)^{2}x^{2}}}{h}$
= $\frac{\frac{x^{2}-x^{2} - 2xh - h^{2}}{(x+h)^{2}x^{2}}}{h}$
= $\frac{\frac{-2xh - h^{2}}{(x+h)^{2}x^{2}}}{h}$
= $\frac{\frac{h(-2x - h)}{(x+h)^{2}x^{2}}}{h}$
= $\frac{-2x - h}{(x+h)^{2}x^{2}}$
\subsection*{答案:}
\begin{align*}
  &= \frac{-2x - h}{(x+h)^{2}x^{2}} \\
\end{align*}

\subsection*{题目: 化简： $3x^{3}-27x$}
\subsection*{思路} 
= $3x^{3} - 3^3x$ 
= $3x(x^{2} - 3^{2})$
= $3x((x-3)*(x+3))$
\subsection*{答案:}
\begin{align*}
  &= 3x((x-3)*(x+3)) \\
\end{align*}

\subsection*{题目: 将下列式子转化为$x^{n}$形式  $\frac{\sqrt{x} \cdot x^{2}}{\sqrt[3]{x}}$}
\subsection*{思路} 
= $\frac{x^{\frac{1}{2}} \cdot x^{2}}{x^{\frac{1}{3}}}$\\
= $\frac{x^{\frac{5}{2}}}{x^{\frac{1}{3}}}$\\
= $x^{\frac{5}{2}-\frac{1}{3}}$\\
= $x^{\frac{15}{6}-\frac{2}{6}}$\\
= $x^{\frac{13}{6}}$\\
\subsection*{答案:}
\begin{align*}
  &= x^{\frac{13}{6}} \\
\end{align*}

\subsection*{题目: 已知 $f(x) = x^{2}-4 $ 计算并化简差商： 
$\frac{f(x+h) - f(x)}{h}$}
\subsection*{思路} 
带入
= $\frac{((x+h)^2-4)-(x^2 -4)}{h}$ 去括号 \\
= $\frac{(x+h)^2-4-x^2 +4}{h}$\\
= $\frac{x^{2} + 2xh + h^{2} - x^{2}}{h}$\\
= $\frac{2xh + h^{2}}{h}$\\
= $\frac{h(2x + h)}{h}$\\
= $2x +h $\\
\subsection*{答案:}
\begin{align*}
  &= 2x +h \\
\end{align*}

\section{Phase 1 Final Exam}
\subsection*{题目: 指数归一化： $\frac{1}{x\sqrt{x}}$}
= $\frac{1}{x \cdot x^{\frac{1}{2}}}$\\
= $\frac{1}{x^{1} \cdot x^{\frac{1}{2}}}$\\
= $\frac{1}{x^{\frac{3}{2}}}$\\
= $x^{-\frac{3}{2}}$\\
\subsection*{答案:}
\begin{align*}
  &= x^{-\frac{3}{2}} \\
\end{align*}

\subsection*{题目: 繁分式化简： $\frac{\frac{2}{x}-\frac{2}{y}}{x-y}$}
= $\frac{\frac{2y-2x}{xy}}{x-y}$\\
= $\frac{\frac{2(y-x)}{xy}}{x-y}$\\
= $\frac{\frac{2 *-(x-y)}{xy}}{x-y}$\\
= $\frac{\frac{-2(x-y)}{xy}}{x-y}$\\
= $\frac{-2}{xy}$\\
= $-2\frac{1}{xy}$\\
\subsection*{答案:}
\begin{align*}
  &= -2\frac{1}{xy} \\
\end{align*}

\subsection*{题目: 因式分解： $3x^{2}-12$}
= $3x^{2} - 3\cdot2^{2}$\\
= $3(x^{2} - 2^{2})$\\
= $3((x+2)(x-2))$\\
\subsection*{答案:}
\begin{align*}
  &= 3((x+2)(x-2)) \\
\end{align*}

\subsection*{题目: 完全平方展开： $(2x - 5)^{2}$}
= $(2x)^{2} - 20x + 5^{2}$\\
= $4x^{2} - 20x + 25$\\
\subsection*{答案:}
\begin{align*}
  &= 4x^{2} - 20x + 25 \\
\end{align*}

\subsection*{题目: 负指数处理： $\begin{pmatrix}
  \frac{2x^{3}}{y^{-2}}
\end{pmatrix}^{-2}$}
= $\frac{1}{\frac{(2x^{3})^{2}}{(y^{-2})^{2}}}$ \\
= $1 \cdot \frac{(y^{-2})^{2}}{(2x^{3})^{2}}$\\
= $\frac{y^{-4}}{2^{2}x^{6}}$\\
= $\frac{\frac{1}{y^{4}}}{4x^{6}}$\\
= $\frac{1}{y^{4}} \cdot \frac{1}{4x^{6}}$\\
= $\frac{1}{y^{4}4x^{6}}$\\
\begin{align*}
  &= \frac{1}{y^{4}4x^{6}} \\
\end{align*}

\subsection*{题目: 二次三项式分解：$x^{2}-x-12$}
=$x^{2} -1x -12$\\
提取 a = -4 , b = 3\\
=$(x - 4)(x + 3)$\\
\begin{align*}
  &= (x - 4)(x + 3) \\
\end{align*}

\subsection*{题目: 立方差商 (The Cube)：$\frac{(x+h)^{3}-x^{3}}{h}$}
= $\frac{((x + h) - x)((x+h)^{2} + x(x+h) + x^{2})}{h}$
= $\frac{h((x+h)^{2} + x(x+h) + x^{2})}{h}$
= $(x+h)^{2} + x(x+h) + x^{2}$
= $x^{2} + 2xh + h^{2} + x^{2} + xh + x^{2}$
= $3x^{2} + 3xh  + h^{2}$


\subsection*{题目: 倒数差商 已知$f(x) = \frac{1}{x}$ (The Reciprocal)：$\frac{f(x+h)-f(x)}{h}$}
= $\frac{\frac{1}{x+h} - \frac{1}{x}}{h}$ \\
= $\frac{\frac{x}{(x+h)x} - \frac{(x+h)}{x(x+h)}}{h}$ \\
= $\frac{\frac{x- (x+h)}{(x+h)x}}{h}$\\
= $\frac{\frac{x- x -h}{(x+h)x}}{h}$\\
= ${-\frac{1}{(x+h)x}}$\\

\begin{align*}
  &= -\frac{1}{(x+h)x} \\
\end{align*}


\subsection*{题目: 根号差商 (The Root)$\frac{\sqrt{x+h}-\sqrt{x}}{h}$}
乘共轭
(a-b)(a+b) = $a^{2}-b^{2}$ \\ 
= $(\sqrt{x+h} - \sqrt{x})(\sqrt{x+h} + \sqrt{x})$ \\  
= $(\sqrt{x+h})^{2} - (\sqrt{x})^{2}$ \\
= $x + h - x$ \\
= $\frac{x + h - x}{(\sqrt{x+h} + \sqrt{x}) \cdot h}$ \\
= $\frac{h}{(\sqrt{x+h} + \sqrt{x}) \cdot h}$ \\
= $\frac{1}{\sqrt{x+h} + \sqrt{x}}$ \\

\begin{align*}
  &=\frac{1}{\sqrt{x+h} + \sqrt{x}} \\
\end{align*}


\subsection*{题目: 连环平方差 $x^{4}-81$}
 = $(x^{2})^{2} - 9^{2}$ \\ 
 = $(x^{2} -9)(x^{2} + 9)$ \\ 
 = $(x^{2} -3^{2})(x^{2} + 3^{2})$ \\ 
 = $((x-3)(x + 3 ))(x^{2} + 3^{2})$\\ 

\begin{align*}
  &= ((x-3)(x + 3 ))(x^{2} + 3^{2}) \\
\end{align*}







\end{document}



